% GENERATED FILE. Edit canonical lessons, then run npm run generate:books.
% canonical-source-hash: fnv1a64:adde799a240ea2c6
% canonical-lessons: ZH-C417-haiguan, ZH-C417-rujing, ZH-C417-chujing, ZH-C417-shenbao, ZH-C417-mianshui

\chapter{Customs}
\label{ch:zh-customs}
\hlchaptermodality{417}

\begin{chapteropening}
\textbf{By the end of this chapter:} \emph{I can go through customs, come into and leave a country, declare things and buy duty-free.}

Complete the last lesson of chapter 417: I can go through customs, come into and leave a country, declare things and buy duty-free.
\end{chapteropening}

\section[hǎiguān]{hǎiguān — customs at the border}

\label{lesson:ZH-C417-haiguan}



\begin{quote}
Before the new one: say the Mandarin for to change planes on the way, then the Mandarin for flying (of an aircraft).
\end{quote}

\subsection*{You'll want to know: hǎiguān}
\emph{hǎiguān} — \textquotedblleft{}customs at the border\textquotedblright{}.

\begin{quote}
\emph{qù hǎiguān zěnme zǒu} — How do I get to customs?
\end{quote}

Its characters come later; for now it is a word for the ear.

\begin{cousinweb}
\emph{hǎiguān} is built from \emph{hǎi} 'the sea' and \emph{guān} 'a barrier', both inherited from Middle Chinese, \emph{xojX} and \emph{kwaen} (Guangyun, in Baxter's transcription), which glosses the second as a bar across a gate; how the pair came to mean customs is not settled.
\end{cousinweb}

\begin{grammarlens}[title={What you've built}]
One more word for getting around.
\end{grammarlens}

\subsection*{Your turn}
\begin{itemize}
\raggedright
  \item \emph{Say it:} \emph{hǎiguān}
  \item \emph{Say it:} \emph{hǎiguān}, once more
  \item \emph{Say it:} \emph{qù hǎiguān zěnme zǒu} — How do I get to customs?
  \item \emph{Recall:} say the Mandarin for to change planes on the way, then the Mandarin for flying (of an aircraft), then say \emph{hǎiguān} again
\end{itemize}

\begin{tcolorbox}[breakable,colback=teal!4,colframe=teal!35!black,title={Before you move on}]
What does hǎiguān mean? (\textquotedblleft{}Customs at the border\textquotedblright{}.) Say it once more.
\end{tcolorbox}
\section[rùjìng]{rùjìng — to enter a country}

\label{lesson:ZH-C417-rujing}



\begin{quote}
Before the new one: say the Mandarin for flying (of an aircraft), then the Mandarin for customs at the border.
\end{quote}

\subsection*{You'll want to know: rùjìng}
\emph{rùjìng} — \textquotedblleft{}to enter a country\textquotedblright{}.

\begin{quote}
\emph{rùjìng yào yòng hùzhào} — You need a passport to enter the country.
\end{quote}

Its characters come later; for now it is a word for the ear.

\begin{cousinweb}
\emph{rùjìng} is built from \zh{入} 'to go in' and \emph{jìng} 'a border', both inherited from Middle Chinese, \emph{nyip} and \emph{kjaengX} (Guangyun, in Baxter's transcription), which glosses the second 'a boundary'.
\end{cousinweb}

\begin{grammarlens}[title={What you've built}]
One more word for getting around.
\end{grammarlens}

\subsection*{Your turn}
\begin{itemize}
\raggedright
  \item \emph{Say it:} \emph{rùjìng}
  \item \emph{Say it:} \emph{rùjìng}, once more
  \item \emph{Say it:} \emph{rùjìng yào yòng hùzhào} — You need a passport to enter the country.
  \item \emph{Recall:} say the Mandarin for flying (of an aircraft), then the Mandarin for customs at the border, then say \emph{rùjìng} again
\end{itemize}

\begin{tcolorbox}[breakable,colback=teal!4,colframe=teal!35!black,title={Before you move on}]
What does rùjìng mean? (\textquotedblleft{}To enter a country\textquotedblright{}.) Say it once more.
\end{tcolorbox}
\section[chūjìng]{chūjìng — to leave a country}

\label{lesson:ZH-C417-chujing}



\begin{quote}
Before the new one: say the Mandarin for customs at the border, then the Mandarin for to enter a country.
\end{quote}

\subsection*{You'll want to know: chūjìng}
\emph{chūjìng} — \textquotedblleft{}to leave a country\textquotedblright{}.

\begin{quote}
\emph{chūjìng de shíhou tāmen kàn le wǒ de hùzhào} — When I left the country, they looked at my passport.
\end{quote}

Its characters come later; for now it is a word for the ear.

\begin{cousinweb}
\emph{chūjìng} is built from \zh{出} 'to go out' and \emph{jìng} 'a border', both inherited from Middle Chinese, \emph{tsyhwit} and \emph{kjaengX} (Guangyun, in Baxter's transcription).
\end{cousinweb}

\begin{grammarlens}[title={What you've built}]
One more word for getting around.
\end{grammarlens}

\subsection*{Your turn}
\begin{itemize}
\raggedright
  \item \emph{Say it:} \emph{chūjìng}
  \item \emph{Say it:} \emph{chūjìng}, once more
  \item \emph{Say it:} \emph{chūjìng de shíhou tāmen kàn le wǒ de hùzhào} — When I left the country, they looked at my passport.
  \item \emph{Recall:} say the Mandarin for customs at the border, then the Mandarin for to enter a country, then say \emph{chūjìng} again
\end{itemize}

\begin{tcolorbox}[breakable,colback=teal!4,colframe=teal!35!black,title={Before you move on}]
What does chūjìng mean? (\textquotedblleft{}To leave a country\textquotedblright{}.) Say it once more.
\end{tcolorbox}
\section[shēnbào]{shēnbào — to declare goods at customs}

\label{lesson:ZH-C417-shenbao}



\begin{quote}
Before the new one: say the Mandarin for to enter a country, then the Mandarin for to leave a country.
\end{quote}

\subsection*{You'll want to know: shēnbào}
\emph{shēnbào} — \textquotedblleft{}to declare goods at customs\textquotedblright{}.

\begin{quote}
\emph{nǐ yǒu dōngxi yào shēnbào ma} — Do you have anything to declare?
\end{quote}

Its characters come later; for now it is a word for the ear.

\begin{cousinweb}
\emph{shēnbào} is built from \emph{shēn} 'to state' and \emph{bào} 'to report', both inherited from Middle Chinese, \emph{syin} and \emph{pawH} (Guangyun, in Baxter's transcription), which glosses the first 'to stretch' and 'to repeat' and the second 'to report'.
\end{cousinweb}

\begin{grammarlens}[title={What you've built}]
One more word for getting around.
\end{grammarlens}

\subsection*{Your turn}
\begin{itemize}
\raggedright
  \item \emph{Say it:} \emph{shēnbào}
  \item \emph{Say it:} \emph{shēnbào}, once more
  \item \emph{Say it:} \emph{nǐ yǒu dōngxi yào shēnbào ma} — Do you have anything to declare?
  \item \emph{Recall:} say the Mandarin for to enter a country, then the Mandarin for to leave a country, then say \emph{shēnbào} again
\end{itemize}

\begin{tcolorbox}[breakable,colback=teal!4,colframe=teal!35!black,title={Before you move on}]
What does shēnbào mean? (\textquotedblleft{}To declare goods at customs\textquotedblright{}.) Say it once more.
\end{tcolorbox}
\section[miǎnshuì]{miǎnshuì — duty-free; tax-free}

\label{lesson:ZH-C417-mianshui}



\begin{quote}
Before the new one: say the Mandarin for to leave a country, then the Mandarin for to declare goods at customs.
\end{quote}

\subsection*{You'll want to know: miǎnshuì}
\emph{miǎnshuì} — \textquotedblleft{}duty-free; tax-free\textquotedblright{}.

\begin{quote}
\emph{jīchǎng lǐ yǒu yì jiā miǎnshuì de shāngdiàn} — There's a duty-free shop in the airport.
\end{quote}

Its characters come later; for now it is a word for the ear.

\begin{cousinweb}
\emph{miǎnshuì} is built from \emph{miǎn} 'to be spared' and \emph{shuì} 'a tax', both inherited from Middle Chinese, \emph{mjenX} and \emph{sywejH} (Guangyun, in Baxter's transcription), which glosses the first 'to put off, to escape' and the second 'a levy'.
\end{cousinweb}

\begin{grammarlens}[title={What you've built}]
One more word for getting around.
\end{grammarlens}

\subsection*{Your turn}
\begin{itemize}
\raggedright
  \item \emph{Say it:} \emph{miǎnshuì}
  \item \emph{Say it:} \emph{miǎnshuì}, once more
  \item \emph{Say it:} \emph{jīchǎng lǐ yǒu yì jiā miǎnshuì de shāngdiàn} — There's a duty-free shop in the airport.
  \item \emph{Recall:} say the Mandarin for to leave a country, then the Mandarin for to declare goods at customs, then say \emph{miǎnshuì} again
\end{itemize}

\begin{tcolorbox}[breakable,colback=teal!4,colframe=teal!35!black,title={Before you move on}]
What does miǎnshuì mean? (\textquotedblleft{}Duty-free; tax-free\textquotedblright{}.) Say it once more.
\end{tcolorbox}
